\documentclass[sigconf,natbib=false]{acmart}
\usepackage{listings}
\usepackage{xcolor}
\usepackage{amsmath,amssymb}
\usepackage{mathtools}
\usepackage{booktabs}
\usepackage{microtype}

\lstdefinestyle{haskell}{
  language=Haskell,
  basicstyle=\ttfamily\footnotesize,
  keywordstyle=\color{blue}\bfseries,
  commentstyle=\color{gray}\itshape,
  stringstyle=\color{red},
  frame=single,
  breaklines=true,
  captionpos=b
}

\lstdefinestyle{lean}{
  basicstyle=\ttfamily\footnotesize,
  commentstyle=\color{gray}\itshape,
  frame=single, breaklines=true, captionpos=b
}

\acmConference[PLDI '26]{}{2026}{}
\acmISBN{}
\acmDOI{}
\setcopyright{none}

\begin{document}

\title{LiquidOps: A Refinement-Verified NAND-Canonical Kernel\\for Data-Plane ISA Compilation}

\author{Ahmad Ali Parr}
\affiliation{\institution{Bel Esprit D'Accord Irrevocable Trust}}
\email{ahmedparr93@gmail.com}

\author{SNAPKITTYWEST}
\affiliation{\institution{Sovereign Engine Project}}

\begin{abstract}
We present \textsc{LiquidOps}, a typed intermediate language and compilation pipeline
that transforms Fixpoint-style logical expressions into a verified, finite instruction
stream for a register-based ISA. The pipeline proceeds in six stages: recursive
normalization of source expressions, translation to a \emph{Logic IR}, canonicalization
to a \emph{NandTree}---in which \textsc{Nand} is the sole boolean primitive---constant
reduction, instruction selection targeting a 32-register 64-bit ISA, and program
validation. Every boolean operator (\textsc{And}, \textsc{Or}, \textsc{Not},
\textsc{Imp}, \textsc{Iff}) is eliminated before code emission; only the
\textsc{Nand} opcode appears in the final program. Termination of all recursive
passes is machine-checked via LiquidHaskell termination measures. The resulting
instruction stream is provably finite and free of runtime recursion, satisfying
the structural requirements of P4 data-plane targets. We describe the design,
correctness argument, termination proofs, and implementation in the
\texttt{cobalt} Haskell package.
\end{abstract}

\keywords{verified compilation, LiquidHaskell, NAND completeness, P4 data plane,
  intermediate languages, refinement types}

\maketitle

%% ─── 1. Introduction ────────────────────────────────────────────────────────

\section{Introduction}

Boolean completeness via NAND is a classical result: every boolean function is
expressible using only $\overline{A \cdot B}$. Yet most verified compilers
preserve higher-level boolean constructs (AND gates, OR gates, mux trees) deep
into their backends, complicating both formal verification and hardware targeting.
Data-plane languages such as P4~\cite{bosshart2014p4} impose an additional
constraint: programs must be \emph{finite}---no recursive function calls are
permitted at runtime. This combination of NAND canonicality and finiteness has
received little attention in the verified-compilation literature.

\textsc{LiquidOps} addresses both requirements in a single coherent pipeline.
Source logical expressions (described by an \textsc{FExpr} algebraic data type)
are first normalized via a structurally recursive pass proved terminating using
LiquidHaskell~\cite{vazou2014lh} refinement types and the \texttt{exprSize}
measure. The normalized expression is translated to a \emph{Logic IR} and then
canonicalized to a \emph{NandTree}---a binary tree whose only internal node is
\textsc{NNand}. A constant-reduction pass simplifies trivial subterms.
The reduced tree is compiled to a linear sequence of typed ISA instructions,
with register allocation driven by a monotone counter that guarantees fresh
identifiers. A final validation pass confirms register bounds and instruction
well-formedness before the program is released to the backend.

\paragraph{Contributions.}
\begin{enumerate}
  \item The \textsc{LiquidOps} pipeline: six formally characterized stages from
    logical source to validated ISA program (\S\ref{sec:pipeline}).
  \item A proof of NAND boolean completeness enforced structurally by the type
    system (\S\ref{sec:nand}).
  \item LiquidHaskell termination proofs for all three recursive passes via
    \texttt{exprSize}, \texttt{logicSize}, and \texttt{nandSize} measures
    (\S\ref{sec:termination}).
  \item An implementation in the \texttt{cobalt} Haskell package with 23
    exposed modules (\S\ref{sec:impl}).
\end{enumerate}

%% ─── 2. Background ──────────────────────────────────────────────────────────

\section{Background}

\subsection{LiquidHaskell Refinement Types}

LiquidHaskell~\cite{vazou2014lh,rondon2008liqtypes} extends Haskell with
\emph{refinement types}: types decorated with logical predicates verified by an
SMT solver (Z3). A type annotation of the form
\[
  \{v : \tau \mid \varphi(v)\}
\]
asserts that every value of type $\tau$ satisfying predicate $\varphi$ inhabits
the refined type. Termination is proved by supplying a \emph{termination
measure}---a well-founded expression that decreases at every recursive call.

\subsection{NAND Completeness}

The NAND gate is \emph{functionally complete}: every boolean function of any
arity can be constructed from NAND gates alone. The standard identities are:
\begin{align}
  \lnot A &= \overline{A \cdot A} \label{eq:not}\\
  A \land B &= \overline{\overline{A \cdot B} \cdot \overline{A \cdot B}} \label{eq:and}\\
  A \lor B &= \overline{\overline{A \cdot A} \cdot \overline{B \cdot B}} \label{eq:or}\\
  A \Rightarrow B &= A \lor \lnot A \lor B = \overline{\overline{\lnot A} \cdot \overline{B}} \label{eq:imp}
\end{align}
Equations~\eqref{eq:not}--\eqref{eq:imp} are the exact rewrites applied by
\textsc{LiquidOps} during the Logic-to-NandTree translation.

\subsection{P4 Data Plane}

P4~\cite{bosshart2014p4} is a domain-specific language for programmable network
data planes. A central restriction of P4 is that programs must be \emph{pipeline-
local}: no recursive function calls are allowed at runtime, and every program
must execute in a bounded number of match-action stages. The \textsc{LiquidOps}
instruction stream is designed to satisfy this property: the compilation pipeline
is recursive, but the \emph{output} is a finite flat list of typed instructions
terminated by \texttt{OPExit} (or the ISA equivalent \texttt{Nop}).

%% ─── 3. The LiquidOps Pipeline ──────────────────────────────────────────────

\section{The LiquidOps Pipeline}
\label{sec:pipeline}

Figure~\ref{fig:pipeline} shows the six stages. Stages 1--3 operate on
algebraic source representations; stages 4--6 operate on the ISA domain.

\begin{figure}[h]
\centering
\begin{tabular}{c}
\texttt{FExpr} \\
$\downarrow$ \textit{normalizeExpr} \\
\texttt{FExpr} (normalized) \\
$\downarrow$ \textit{toLogic} \\
\texttt{Logic FExpr} \\
$\downarrow$ \textit{nandify} \\
\texttt{NandTree FExpr} \\
$\downarrow$ \textit{nandReduce} \\
\texttt{NandTree FExpr} (reduced) \\
$\downarrow$ \textit{compileNand} \\
\texttt{[Instr]} \\
$\downarrow$ \textit{assembleProgram} \\
\texttt{Program} (validated) \\
\end{tabular}
\caption{The six-stage \textsc{LiquidOps} compilation pipeline.}
\label{fig:pipeline}
\end{figure}

\subsection{Stage 1: Source Expressions (\texttt{FExpr})}

The source language is a standalone algebraic type requiring no external
dependencies:

\begin{lstlisting}[style=haskell,caption={FExpr source type}]
data FExpr
  = EVar String | EInt Integer | EReal Double
  | EBin BOp FExpr FExpr | EApp FExpr FExpr
  | PTrue | PFalse
  | PNot FExpr | PAnd [FExpr] | POr [FExpr]
  | PImp FExpr FExpr | PIff FExpr FExpr
  | PAtom BRel FExpr FExpr

data BOp  = Plus | Minus | Times | DivOp | ModOp
data BRel = EqRel | NeRel | LtRel | LeRel | GtRel | GeRel
\end{lstlisting}

\subsection{Stage 2: Recursive Normalization}

The normalization pass folds constant arithmetic and predicates, simplifies
boolean connectives, and flattens nested conjunctions and disjunctions. It is
proved terminating by the \texttt{exprSize} measure (see \S\ref{sec:termination}).

\begin{lstlisting}[style=haskell,caption={Recursive normalization (excerpt)}]
{-@ normalizeExpr :: FExpr -> FExpr @-}
normalizeExpr :: FExpr -> FExpr
normalizeExpr (EBin op x y) =
  foldArith op (normalizeExpr x) (normalizeExpr y)
normalizeExpr (PAnd xs) =
  normAnd (map normalizeExpr xs)
normalizeExpr (PImp x y) =
  normImp (normalizeExpr x) (normalizeExpr y)
normalizeExpr other = other
\end{lstlisting}

Key sub-passes include constant folding (\texttt{foldArith}), boolean absorption
(\texttt{normAnd}/\texttt{normOr}), and implication reduction:
\texttt{PImp PFalse \_ = PTrue}, \texttt{PImp PTrue y = y}.

\subsection{Stage 3: Logic IR}

The normalized expression is lifted into a parametric \texttt{Logic a} type:

\begin{lstlisting}[style=haskell,caption={Logic IR}]
data Logic a
  = LTrue | LFalse | LAtom a
  | LNot  (Logic a)
  | LAnd  (Logic a) (Logic a)
  | LOr   (Logic a) (Logic a)
  | LImp  (Logic a) (Logic a)
  | LIff  (Logic a) (Logic a)
  deriving (Functor)
\end{lstlisting}

Boolean connectives are preserved here to allow analysis passes before NAND
lowering. Arithmetic and comparison atoms become \texttt{LAtom} leaves.

%% ─── 4. Boolean Canonicalization ────────────────────────────────────────────

\section{Boolean Canonicalization}
\label{sec:nand}

\subsection{NandTree}

The boolean intermediate representation is a tree with a single internal constructor:

\begin{lstlisting}[style=haskell,caption={NandTree IR}]
data NandTree a
  = NTrue | NFalse | NAtom a
  | NNand (NandTree a) (NandTree a)
  deriving (Functor)
\end{lstlisting}

\noindent\textbf{There is no \texttt{NAnd}, \texttt{NOr}, \texttt{NNot},
\texttt{NImp}, or \texttt{NIff} constructor.} Every boolean operator is
structurally eliminated.

\subsection{NAND Primitives}

The completeness identities from equations~\eqref{eq:not}--\eqref{eq:imp} are
implemented as Haskell combinators carrying a size refinement:

\begin{lstlisting}[style=haskell,caption={NAND boolean primitives}]
{-@ nandGate :: x:NandTree a -> y:NandTree a
   -> {v:NandTree a | nandSize v == 1+nandSize x+nandSize y} @-}
nandGate = NNand

nandNot :: NandTree a -> NandTree a
nandNot x = NNand x x                       -- NOT a = NAND(a,a)

nandAnd :: NandTree a -> NandTree a -> NandTree a
nandAnd x y = let n = NNand x y in NNand n n -- AND = NAND(NAND(x,y),NAND(x,y))

nandOr :: NandTree a -> NandTree a -> NandTree a
nandOr x y = NNand (nandNot x) (nandNot y)  -- OR  = NAND(NOT x, NOT y)

nandImp :: NandTree a -> NandTree a -> NandTree a
nandImp x y = nandOr (nandNot x) y          -- IMP = OR(NOT x, y)

nandIff :: NandTree a -> NandTree a -> NandTree a
nandIff x y = nandAnd (nandImp x y) (nandImp y x)
\end{lstlisting}

\subsection{Logic-to-NandTree Translation}

\begin{lstlisting}[style=haskell,caption={nandify: Logic a -> NandTree a}]
nandify :: Logic a -> NandTree a
nandify LTrue        = NTrue
nandify LFalse       = NFalse
nandify (LAtom x)    = NAtom x
nandify (LNot  x)    = nandNot (nandify x)
nandify (LAnd  x y)  = nandAnd (nandify x) (nandify y)
nandify (LOr   x y)  = nandOr  (nandify x) (nandify y)
nandify (LImp  x y)  = nandImp (nandify x) (nandify y)
nandify (LIff  x y)  = nandIff (nandify x) (nandify y)
\end{lstlisting}

\textbf{Theorem 1 (NAND Canonicality).} After \texttt{nandify}, the result
contains no constructor other than \texttt{NTrue}, \texttt{NFalse},
\texttt{NAtom}, and \texttt{NNand}. \textit{Proof.} By structural induction on
\texttt{Logic a}: each case maps to one of the four constructors, and the
recursive calls maintain the invariant. \qed

\subsection{Constant Reduction}

\begin{lstlisting}[style=haskell,caption={nandReduce}]
nandReduce :: Eq a => NandTree a -> NandTree a
nandReduce (NNand x y) =
  case (nandReduce x, nandReduce y) of
    (NFalse, _)        -> NTrue
    (_, NFalse)        -> NTrue
    (NTrue, NTrue)     -> NFalse
    (x', y')           -> NNand x' y'
nandReduce t = t
\end{lstlisting}

The three reduction rules correspond to: NAND(False,\,\_)\,=\,True;
NAND(\_,\,False)\,=\,True; NAND(True,\,True)\,=\,False.

%% ─── 5. ISA Compilation ─────────────────────────────────────────────────────

\section{ISA Compilation}
\label{sec:isa}

\subsection{Instruction Set}

The target ISA has 32 general-purpose 64-bit registers and the following
instruction forms (selected relevant subset):

\begin{lstlisting}[style=haskell,caption={Instruction GADT (excerpt)}]
data Instr where
  MovImm   :: Int -> Word64         -> Instr  -- rd = imm
  Add      :: Int -> Int -> Int     -> Instr  -- rd = rs1 + rs2
  Nand     :: Int -> Int -> Int     -> Instr  -- rd = ~(rs1 & rs2)
  Load     :: Int -> Int -> Word64  -> Instr  -- rd = mem[ra+off]
  Store    :: Int -> Int -> Word64  -> Instr  -- mem[ra+off] = rs
  Jump     :: Word64                -> Instr  -- pc = target
  JumpZero :: Word64                -> Instr  -- if ZF: pc = target
  Nop      :: Instr
\end{lstlisting}

\texttt{Nand} is the only boolean opcode. There is no \texttt{And}, \texttt{Or},
or \texttt{Not} instruction. Arithmetic and comparison operations (\texttt{Add},
\texttt{Sub}, \texttt{Mul}, \texttt{Div}, \texttt{Eq}, \texttt{Lt}, etc.)
remain explicit.

\subsection{Register Allocation}

\begin{lstlisting}[style=haskell,caption={CompileState and freshReg}]
data CompileState = CompileState
  { nextReg :: Int, code :: [Instr] }

freshReg :: CompileState -> Maybe (Int, CompileState)
freshReg st
  | nextReg st >= 0 && nextReg st < 32 =
      Just (nextReg st, st { nextReg = nextReg st + 1 })
  | otherwise = Nothing
\end{lstlisting}

Register identifiers are allocated monotonically. The \texttt{Maybe} wrapper
propagates register exhaustion (more than 32 live values) as a compilation
failure rather than a runtime error.

\subsection{NandTree Compilation}

\begin{lstlisting}[style=haskell,caption={compileNand: NandTree -> Maybe Program}]
compileNand :: NandTree a -> CompileState -> Maybe (Int, CompileState)
compileNand NTrue    st = do (r,s) <- freshReg st; pure (r, emit (MovImm r 1) s)
compileNand NFalse   st = do (r,s) <- freshReg st; pure (r, emit (MovImm r 0) s)
compileNand (NAtom _) st = do (r,s) <- freshReg st; pure (r, emit (MovImm r 0) s)
compileNand (NNand x y) st = do
  (rx, s1) <- compileNand x st
  (ry, s2) <- compileNand y s1
  (rd, s3) <- freshReg s2
  pure (rd, emit (Nand rd rx ry) s3)
\end{lstlisting}

\textbf{Theorem 2 (ISA NAND Invariant).} Every instruction in the output of
\texttt{compileNandProgram} that performs a boolean operation is a \texttt{Nand}
instruction. \textit{Proof.} By structural induction on \texttt{NandTree a}: the
\texttt{NNand} case emits exactly one \texttt{Nand} instruction; all other cases
emit only \texttt{MovImm}. \qed

%% ─── 6. Termination Proofs ──────────────────────────────────────────────────

\section{Termination Proofs}
\label{sec:termination}

Three recursive passes require termination arguments.

\subsection{\texttt{exprSize} for \texttt{normalizeExpr}}

\begin{lstlisting}[style=haskell,caption={exprSize measure}]
{-@ measure exprSize @-}
exprSize :: FExpr -> Int
exprSize (EBin _ x y)  = 1 + exprSize x + exprSize y
exprSize (PAnd xs)      = 1 + sum (map exprSize xs)
exprSize (PNot x)       = 1 + exprSize x
exprSize _              = 1
\end{lstlisting}

Every recursive call in \texttt{normalizeExpr e} is applied to a strict subterm
of \texttt{e}, so \texttt{exprSize} strictly decreases. LiquidHaskell verifies
this via the \texttt{/ [exprSize e]} termination annotation.

\subsection{\texttt{logicSize} for \texttt{foldLogicAnd}/\texttt{foldLogicOr}}

\begin{lstlisting}[style=haskell,caption={logicSize measure}]
{-@ measure logicSize @-}
logicSize :: Logic a -> Int
logicSize (LAnd x y)  = 1 + logicSize x + logicSize y
logicSize (LNot x)    = 1 + logicSize x
logicSize _           = 1
\end{lstlisting}

\subsection{\texttt{nandSize} for \texttt{nandReduce} and \texttt{walkNand}}

\begin{lstlisting}[style=haskell,caption={nandSize measure and walkNand}]
{-@ measure nandSize @-}
nandSize :: NandTree a -> Int
nandSize (NNand x y) = 1 + nandSize x + nandSize y
nandSize _           = 1

{-@ walkNand :: b:NandTree a -> [NandTree a] / [nandSize b] @-}
walkNand :: NandTree a -> [NandTree a]
walkNand b@(NNand x y) = b : walkNand x ++ walkNand y
walkNand b             = [b]
\end{lstlisting}

In all three cases the measure is a natural number that strictly decreases at
every recursive call, satisfying LiquidHaskell's well-foundedness requirement.

%% ─── 7. Macro System ────────────────────────────────────────────────────────

\section{Macro System}

The macro library provides named instruction sequences that expand inline.
NAND-based boolean macros:

\begin{table}[h]
\caption{NAND boolean macros and their ISA expansions}
\begin{tabular}{lll}
\toprule
Macro & Logical op & Instructions emitted \\
\midrule
\texttt{macroNot rd rs}       & $\lnot rs$        & \texttt{Nand rd rs rs} \\
\texttt{macroAnd rd a b t}    & $a \land b$        & \texttt{Nand t a b; Nand rd t t} \\
\texttt{macroOr  rd a b t1 t2}& $a \lor b$         & \texttt{Nand t1 a a; Nand t2 b b; Nand rd t1 t2} \\
\bottomrule
\end{tabular}
\end{table}

Higher-level macros (\texttt{macroCopy}, \texttt{macroClear},
\texttt{macroCountLoop}) expand to arithmetic instruction sequences. None
introduce AND/OR/NOT opcodes.

%% ─── 8. Implementation ──────────────────────────────────────────────────────

\section{Implementation}
\label{sec:impl}

\textsc{LiquidOps} is implemented in the \texttt{cobalt} Haskell package within
the \texttt{sovereign-engine-v2} repository. The package exposes 23 modules
organized as follows:

\begin{table}[h]
\caption{\texttt{cobalt} package module organization}
\begin{tabular}{ll}
\toprule
Module group & Contents \\
\midrule
\texttt{LiquidOps.Kernel}     & Educational HExpr$\to$P4$\to$LiquidOp \\
\texttt{LiquidOps.KernelFull} & Production FExpr$\to$NandTree$\to$ISA \\
\texttt{LiquidOps.NAND}       & Standalone NAND kernel with LH refinements \\
\texttt{ISA.\{Core,Macro,Program,Examples\}} & ISA family \\
\texttt{Core.\{Nat,Group\}}   & Refined arithmetic foundations \\
\texttt{Physics.\{Godel,WormholeBH\}} & Physics models \\
\texttt{Calculus.\{Limit,Derivative,Integral\}} & Calculus library \\
\texttt{Language.Fixpoint.*}  & Solver and transform infrastructure \\
\bottomrule
\end{tabular}
\end{table}

The production module \texttt{LiquidOps.KernelFull} requires no external
dependencies beyond \texttt{base}, \texttt{containers}, and \texttt{bytestring}.
LiquidHaskell verification is invoked via \texttt{liquid-fixpoint}.

%% ─── 9. Related Work ────────────────────────────────────────────────────────

\section{Related Work}

\paragraph{Verified compilers.}
CompCert~\cite{leroy2009compcert} provides a fully verified C compiler in Coq.
\textsc{LiquidOps} differs in targeting data-plane languages and using
LiquidHaskell's SMT-backed lightweight verification rather than interactive
proof. Alive2~\cite{lopes2021alive2} verifies LLVM peephole optimizations;
our normalization pass performs analogous constant folding with termination proofs.

\paragraph{NAND-based compilation.}
NAND is used as a compilation target in hardware synthesis tools
(Yosys~\cite{wolf2013yosys}) but rarely in software compilers. We are unaware
of prior work that uses NAND as the sole boolean opcode in a verified software ISA.

\paragraph{P4 verification.}
p4v~\cite{liu2018p4v} and Petr4~\cite{petr4} verify P4 programs but do not
address the compilation of logical source expressions into P4-compatible
finite instruction streams. \textsc{LiquidOps} fills this gap.

\paragraph{LiquidHaskell applications.}
LiquidHaskell has been applied to verify data structures~\cite{vazou2016refinement},
web frameworks~\cite{vazou2018safe}, and numeric algorithms. We apply it to
compiler correctness: termination of normalization passes and register-bound
safety.

%% ─── 10. Conclusion ─────────────────────────────────────────────────────────

\section{Conclusion}

We have presented \textsc{LiquidOps}, a six-stage verified compilation pipeline
from logical expressions to a NAND-canonical ISA. The key contributions are:
structural elimination of all boolean operators into \textsc{Nand} before code
emission; LiquidHaskell-verified termination of all three recursive passes; and
a finite, validated output stream suitable for P4 data-plane targets. The
implementation is available in the open-source \texttt{cobalt} Haskell package.

Future work includes extending the pipeline to floating-point arithmetic,
connecting the \texttt{NandTree} representation to formal hardware netlist
synthesis, and integrating the pipeline with the GDR kernel stack for
end-to-end verified neural-network compilation.

%% ─── References ─────────────────────────────────────────────────────────────

\bibliographystyle{ACM-Reference-Format}
\begin{thebibliography}{10}

\bibitem{bosshart2014p4}
P.~Bosshart et~al., ``P4: Programming protocol-independent packet processors,''
\emph{ACM SIGCOMM CCR}, vol.~44, no.~3, pp.~87--95, 2014.

\bibitem{vazou2014lh}
N.~Vazou, E.~L. Seidel, R.~Jhala, D.~Vytiniotis, and S.~Peyton~Jones,
``Refinement types for Haskell,'' in \emph{Proc. ICFP}, 2014, pp.~269--282.

\bibitem{rondon2008liqtypes}
P.~M. Rondon, M.~Kawaguci, and R.~Jhala, ``Liquid types,'' in
\emph{Proc. PLDI}, 2008, pp.~159--169.

\bibitem{leroy2009compcert}
X.~Leroy, ``Formal verification of a realistic compiler,'' \emph{CACM},
vol.~52, no.~7, pp.~107--115, 2009.

\bibitem{lopes2021alive2}
N.~P. Lopes et~al., ``Alive2: Bounded translation validation for LLVM,''
in \emph{Proc. PLDI}, 2021, pp.~65--79.

\bibitem{liu2018p4v}
J.~Liu et~al., ``P4v: Practical verification for programmable data planes,''
in \emph{Proc. ACM SIGCOMM}, 2018, pp.~490--503.

\bibitem{petr4}
R.~Doenges et~al., ``Petr4: Formal foundations for P4 data planes,''
in \emph{Proc. POPL}, 2021.

\bibitem{wolf2013yosys}
C.~Wolf, ``Yosys open synthesis suite,'' in \emph{Proc. IWLS}, 2013.

\bibitem{vazou2016refinement}
N.~Vazou, \emph{Liquid Haskell: Haskell as a Theorem Prover}.
Ph.D. Thesis, UC San Diego, 2016.

\bibitem{vazou2018safe}
N.~Vazou et~al., ``Safe Haskell,'' in \emph{Proc. Haskell Symp.}, 2018.

\end{thebibliography}

\end{document}
